\Askhsh{21636}{2}
\begin{ylh}{1}
    \item 9.2. Το Πυθαγόρειο θεώρημα
    \item 10.5. Λόγος εμβαδών όμοιων τριγώνων - πολυγώνων.
\end{ylh}
Δίνεται τρίγωνο $\text{AB\Gamma}$, με μήκη πλευρών $\text{AB}=6$, $\text{A\Gamma}=8$, και $\text{B\Gamma}=10$.
\begin{alist}
\item Να αποδείξετε ότι το τρίγωνο $\text{AB\Gamma}$ είναι ορθογώνιο. \monades{10}
\item Αν $\text{A\Delta}$ είναι ύψος του τριγώνου $\text{AB\Gamma}$ τότε:
    \begin{rlist}
    \item Να αποδείξετε ότι τα τρίγωνα $\text{AB\Delta}$ και $\text{A\Gamma\Delta}$ είναι όμοια με λόγο ομοιότητας $\lambda = \dfrac{3}{4}$. \monades{10}
    \item Να υπολογίσετε το λόγο: $\dfrac{(\text{AB\Delta})}{(\text{A\Gamma\Delta})}$. \monades{5}
    \end{rlist}
\end{alist}